\chapter{Coronary Artery Bypass Graft Surgery}

\section*{Overview}
Coronary artery bypass grafting (CABG) creates new routes for blood to flow around severely narrowed or blocked coronary arteries.

\section*{Why This Test or Procedure Is Ordered}
It may be recommended for left-main or multivessel coronary disease, diabetes with extensive disease, symptoms despite treatment, or anatomy unsuitable for angioplasty.

\section*{How This Test or Procedure Is Performed}
Under general anesthesia, the surgeon uses segments of the internal mammary artery, radial artery, or vein to bypass blocked coronary segments. The operation may use a heart-lung machine or a beating-heart technique.

\section*{Risks and Recovery}
Risks include bleeding, infection, stroke, heart attack, rhythm problems, kidney injury, graft failure, and death. Recovery involves intensive monitoring, gradual activity, medicines, risk-factor control, and cardiac rehabilitation.

\section*{References}
\begin{enumerate}
  \item MedlinePlus. Coronary Artery Bypass Surgery. U.S. National Library of Medicine; 2025.
  \item Current Medical Diagnosis and Treatment. McGraw-Hill; 2025.
\end{enumerate}
\clearpage
